\section*{Capítulo \ref{ch:g11:lithosphere-asthenosphere} --- Litosfera e astenosfera}

\begin{solution}{exo:g11:lithosphere-asthenosphere:1}
A \omterm{def:g11:lithosphere-asthenosphere:lithosphere}{litosfera} é a camada externa rígida da Terra, feita da crosta e do manto
mais alto; a \omterm{def:g11:lithosphere-asthenosphere:lithosphere}{astenosfera} é o manto mais quente e mais fraco abaixo dela, que
escoa lentamente ao longo de grandes intervalos de tempo. A crosta oceânica,
a crosta continental e o manto mais alto fazem parte da \omterm{def:g11:lithosphere-asthenosphere:lithosphere}{litosfera}; o núcleo
externo, não.
\end{solution}

\begin{solution}{exo:g11:lithosphere-asthenosphere:2}
O iceberg afunda até deslocar o seu próprio peso de água: a fração debaixo
da água é $910/1025 \approx 0.888$, cerca de \qty{89}{\percent}. A fração
acima é de cerca de \qty{11.2}{\percent}, logo a altura total é
$20/0.112 \approx \qty{180}{m}$, dos quais cerca de \qty{160}{m} ficam
debaixo da água.
\end{solution}

\begin{solution}{exo:g11:lithosphere-asthenosphere:3}
$15 + 26 \times 3 = \qty{93}{\celsius}$. A \qty{100}{km}, o mesmo gradiente
daria mais de \qty{2600}{\celsius}, muito acima da temperatura de fusão da
rocha do manto; mas as \omterm{def:g11:seismic-structure:body-waves}{ondas S} atravessam o manto superior, que portanto é
sólido. O gradiente precisa ficar muito menor com a profundidade.
\end{solution}

\begin{solution}{exo:g11:lithosphere-asthenosphere:4}
A faixa de 0 a \qty{500}{m} acima do nível do mar, com cerca de
\qty{16}{\percent} da superfície, e a faixa de 4000 a \qty{4500}{m} abaixo
do nível do mar, com cerca de \qty{11}{\percent} (a faixa de 4500 a
\qty{5000}{m} abaixo é quase do mesmo tamanho).
\end{solution}

\begin{solution}{exo:g11:lithosphere-asthenosphere:5}
Para $P = 0$: \qty{1121}{\celsius}. A \qty{3}{GPa}: $-5.104 \times 9 +
132.899 \times 3 + 1120.661 \approx -45.9 + 398.7 + 1120.7 \approx
\qty{1473}{\celsius}$.
\end{solution}

\begin{solution}{exo:g11:lithosphere-asthenosphere:6}
Cerca de \qty{730}{\celsius} a \qty{50}{km} (\omterm{def:g11:lithosphere-asthenosphere:lithosphere}{litosfera}), cerca de
\qty{1370}{\celsius} a \qty{150}{km} e \qty{1410}{\celsius} a
\qty{250}{km} (\omterm{def:g11:lithosphere-asthenosphere:lithosphere}{astenosfera}). A temperatura mais sobe entre 50 e
\qty{150}{km}, em mais de \qty{600}{\celsius}, porque a maior parte desse
intervalo fica na \omterm{def:g11:lithosphere-asthenosphere:lithosphere}{litosfera}; entre 150 e \qty{250}{km}, ela sobe apenas
cerca de \qty{45}{\celsius}.
\end{solution}

\begin{solution}{exo:g11:lithosphere-asthenosphere:7}
Raiz: $3 \times 2830/550 \approx \qty{15.4}{km}$. Espessura a mais:
$3 \times 3380/550 \approx \qty{18.4}{km}$, isto é, os \qty{3}{km} da
cadeia mais a sua raiz.
\end{solution}

\begin{solution}{exo:g11:lithosphere-asthenosphere:8}
A crosta é mais espessa em $55 - 41.1 = \qty{13.9}{km}$, e $\Delta H = h
\times 3380/550$, logo $h = 13.9 \times 550/3380 \approx \qty{2.3}{km}$.
\end{solution}

\begin{solution}{exo:g11:lithosphere-asthenosphere:9}
Diminuição: $(4.470 - 4.419)/4.470 \approx \qty{1.1}{\percent}$. Salto:
$(4.644 - 4.419)/4.419 \approx \qty{5.1}{\percent}$. Com a profundidade, a
rocha fica cada vez mais comprimida, o que torna as ondas sísmicas mais
rápidas; uma diminuição significa que outra coisa prevalece: a rocha está
quente e perto da fusão, e mais fraca.
\end{solution}

\begin{solution}{exo:g11:lithosphere-asthenosphere:10}
(1) $(1345 - 610)/60 \approx \qty{12}{\celsius}$ por quilômetro. (2) É cerca
da metade do gradiente na superfície. Parte do calor é produzida dentro da
crosta superior; esse calor também precisa atravessar as camadas de cima
para escapar, o que torna o gradiente mais forte perto da superfície. Abaixo
da crosta, apenas o calor que vem de mais fundo atravessa a rocha, e o
gradiente é menor.
\end{solution}

\begin{solution}{exo:g11:lithosphere-asthenosphere:11}
A \qty{200}{km}: $1300 + 90 = \qty{1390}{\celsius}$; a \qty{400}{km}:
$1300 + 180 = \qty{1480}{\celsius}$. O aumento é de \qty{90}{\celsius} em
\qty{200}{km}. Através da \omterm{def:g11:lithosphere-asthenosphere:lithosphere}{litosfera}, a temperatura sobe
$1345 - 15 = \qty{1330}{\celsius}$ em apenas \qty{100}{km}: cerca de trinta
vezes mais por quilômetro.
\end{solution}

\begin{solution}{exo:g11:lithosphere-asthenosphere:12}
(1) $118 + 12.3\,z = 1300 + 0.45\,z$ dá $z = 1182/11.85 \approx
\qty{100}{km}$. (2) $118 + 21.5\,z = 1300 + 0.45\,z$ dá $z = 1182/21.05
\approx \qty{56}{km}$. (3) O oeste, mais quente, tem uma \omterm{def:g11:lithosphere-asthenosphere:lithosphere}{litosfera} muito
mais fina, como foi medido (cerca de \qty{40}{km} onde ela é mais fina); a
\omterm{def:g11:lithosphere-asthenosphere:lithosphere}{litosfera} muito espessa do meio do continente precisa de um gradiente ainda
menor.
\end{solution}

\begin{solution}{exo:g11:lithosphere-asthenosphere:13}
(1) Pela proposição, $\Delta H = h\,\rho_m/(\rho_m - \rho_c)$, logo $h =
\Delta H\,(\rho_m - \rho_c)/\rho_m$. (2) $550/3380 \approx 0.16$: cerca de um
sexto de qualquer espessura a mais de crosta aparece como altura, cinco
sextos como raiz. (3) Para o iceberg, a fração acima da água é
$(1025 - 910)/1025 \approx 0.11$: a mesma regra, (densidade do fluido $-$
densidade do corpo que flutua)/densidade do fluido.
\end{solution}

\begin{solution}{exo:g11:lithosphere-asthenosphere:14}
(1) $1489 - 1345 = \qty{144}{\celsius}$. (2) Um pouco de material fundido se
junta em películas finas entre os grãos da rocha. Os grãos então deixam de
estar firmemente presos uns aos outros: a rocha resiste menos ao
cisalhamento, de modo que as \omterm{def:g11:seismic-structure:body-waves}{ondas S} ficam mais lentas, e ela pode se
deformar e escoar com muito mais facilidade.
\end{solution}

\begin{solution}{exo:g11:lithosphere-asthenosphere:15}
(1) Não. Aplainar as terras deixaria a crosta continental um pouco mais
fina, mas ela continuaria espessa e leve, e continuaria flutuando muito mais
alto que a crosta fina e densa do fundo do oceano; os oceanos continuariam
enchendo as bacias entre os continentes. (2) Existem dois tipos de crosta,
de espessura e densidade diferentes: uma crosta continental espessa e leve e
uma crosta oceânica fina e densa.
\end{solution}

\begin{solution}{pb:g11:lithosphere-asthenosphere:1}
\textbf{1.} $5.01 - 0.67 = \qty{4.34}{km}$.

\textbf{2.} $(3380 - 2830)/3380 \approx \qty{16}{\percent}$.

\textbf{3.} A fração $2830/3380 \approx 0.84$ ficaria abaixo; o iceberg
afunda um pouco mais, $910/1025 \approx 0.89$, porque o gelo tem uma
densidade mais próxima da água do que a crosta da do manto.

\textbf{4.} Durante os segundos de uma onda sísmica, ou o tempo de um
terremoto, o manto é um sólido; só ao longo de intervalos de tempo muito
longos a \omterm{def:g11:lithosphere-asthenosphere:lithosphere}{astenosfera} escoa e deixa a crosta flutuar.

\textbf{5.} Abaixo do \omterm{def:g11:seismic-structure:moho}{Moho} mais profundo, as duas colunas são feitas da
mesma rocha do manto nas mesmas profundidades, e ali pesam o mesmo; todas as
diferenças ficam acima dessa profundidade e precisam ser incluídas.

\textbf{6.} O continente médio: crosta de \qty{0.67}{km} acima a
\qty{40.4}{km} abaixo do nível do mar, depois manto até \qty{70}{km}. O
planalto: crosta de \qty{5.01}{km} acima do nível do mar até um \omterm{def:g11:seismic-structure:moho}{Moho} mais
fundo pela raiz, depois manto até \qty{70}{km}.

\textbf{7.} A coluna do planalto tem $h + r$ de crosta a mais, e a coluna
do continente tem, em vez disso, manto na espessura $r$: $\rho_c\,(h + r) = \rho_m\,r$,
logo $r = h\,\rho_c/(\rho_m - \rho_c)$ e $\Delta H = h + r =
h\,\rho_m/(\rho_m - \rho_c)$.

\textbf{8.} $3380/550 \approx 6.15$.

\textbf{9.} $\Delta H = 4.34 \times 6.15 \approx \qty{26.7}{km}$; a raiz é
$26.7 - 4.34 \approx \qty{22.3}{km}$.

\textbf{10.} O \omterm{def:g11:seismic-structure:moho}{Moho} do continente médio fica $41.1 - 0.67 \approx
\qty{40.4}{km}$ abaixo do nível do mar; o do planalto fica $40.4 + 22.3 \approx
\qty{62.8}{km}$ abaixo do nível do mar.

\textbf{11.} $41.1 + 26.7 \approx \qty{67.8}{km}$.

\textbf{12.} Um quilômetro de rocha acima de um metro quadrado tem uma massa
de $\rho \times \qty{1000}{m^3}$: \num{2830} toneladas de crosta ou
\num{3380} toneladas de manto. Continente: $41.1 \times 2830 + (70 - 40.43) \times 3380
\approx \num{116300} + \num{99900} \approx \num{216300}$ toneladas. Planalto:
$67.77 \times 2830 + (70 - 62.76) \times 3380 \approx \num{191800} +
\num{24500} \approx \num{216300}$ toneladas. As duas são iguais.

\textbf{13.} $41.1 - 0.67 \times 6.15 \approx \qty{37.0}{km}$.

\textbf{14.} $67.8/41.1 \approx 1.65$ vez.

\textbf{15.} $3300/470 \approx 7.02$, $\Delta H \approx 4.34 \times 7.02
\approx \qty{30.5}{km}$, e uma crosta de cerca de \qty{71.6}{km}. Uma
mudança de cerca de \qty{2}{\percent} na densidade do manto muda a espessura
em cerca de \qty{4}{km}: o resultado depende visivelmente das densidades,
mas continua perto de \qty{70}{km}.

\textbf{16.} O modelo dá cerca de \qty{68}{km}, as ondas sísmicas cerca de
\qty{70}{km}: eles concordam dentro de poucos por cento.

\textbf{17.} $h = 3.23 - 0.67 = \qty{2.56}{km}$, $\Delta H = 2.56 \times
6.15 \approx \qty{15.7}{km}$, e uma crosta de cerca de \qty{56.8}{km}, uns
\qty{7}{km} a mais que o medido.

\textbf{18.} Mais fina. Uma coluna cuja parte de cima é mais leve que a
média pesa menos para a mesma espessura, logo flutua mais alto: uma crosta
mais leve precisa de uma raiz menor para ficar a uma altura dada. A crosta
real pode portanto ser mais fina que a do modelo, como foi medido.

\textbf{19.} Ao longo de grandes intervalos de tempo, a \omterm{def:g11:lithosphere-asthenosphere:lithosphere}{astenosfera} escoa:
ela se afasta sob uma coluna pesada e volta sob uma coluna leve, até que as
colunas se equilibrem.

\textbf{20.} \textbf{A crosta sob o Tibete tem cerca de \qty{70}{km} de
espessura}, muito mais que os \qty{41}{km} de um continente médio: o
planalto fica a cinco quilômetros de altura porque a sua crosta espessa e
leve flutua sobre o manto, mais denso, com uma raiz de mais de \qty{20}{km}.
\end{solution}
